\documentclass{report}
\usepackage{amsmath,amssymb}
\setlength{\parindent}{0mm}
\setlength{\parskip}{1em}
\begin{document}
\begin{center}
\rule{6in}{1pt} \
{\large Hao LIU \\
{\bf Local Defect Correction multigrid Method coupled with the Zienkiewicz-Zhu \emph{a posteriori} error estimator in elastostatics Solid Mechanics}}

CEA Cadarache \\ DEN/DEC/SESC/LSC \\ Building 151 \\ 13108 SAINT-PAUL-LEZ-DURANCE CEDEX \\ France
\\
{\tt hao.liu@cea.fr}\\
Isabelle Rami{\`e}re\\
Fr{\'e}d{\'e}ric Lebon\end{center}

Local multigrid methods provide a technique to solve problems which
locally require high accuracy (cracks, discontinuous boundary conditions,
entrant corners etc.) in acceptable computational times and memory space.
These methods have several advantages such as fast resolution on each
sub-grid, work on local fine meshes only...

We propose here to use the Local Defect Correction (LDC) multigrid
method, proposed by Hackbush \cite{Hackbusch84}, to solve Solid Mechanics
Problems. From an initial coarse mesh, this method focuses on recursively
adding local sub-grids in areas where higher accuracy is required. We
keep adding levels of refinement until the finest grid has the desired
accuracy. Then prolongation and restriction operators are used to link
the different level of grids. It should be noted that this method is
generic : solver, refinement ratio , mesh type, model etc. could be
different (or not) between the several levels of refinement.

In order to automatically detect the zones to be refined, we propose to
couple the LDC multigrid method with the Zienkiewicz-Zhu \emph{a
posteriori} error estimator \cite{Zienkiewicz87}. Based on the fact that
the stress field calculated by the finite element method is discontinuous
between elements, the Zienkiewicz-Zhu \emph{a posteriori} error estimator
constructs a smoothed stress field. Then the error estimator is evaluated
by the difference between the finite element stress solution and the
smoothed stress fields. A strategy of coupling the LDC multigrid method
with the Zienkiewicz-Zhu \emph{a posteriori} error was already introduced
in \cite{Barbie13CAS,Barbie11}. The authors proposed to refine the
elements for which the indicator is greater than $\alpha$\% of the
maximum error. This strategy is easy to implement but requires the
knowledge of the $\alpha$ (which may depend of the problem) and the
maximal number of sub-grids (because this indicator never stops). We
propose here an automatic procedure of refinement working directly on the
elements for which the stress error is superior than a threshold
(relative error) given by the user. In this way, we avoid the dependence
on a arbitrary coefficient.

We apply this method on an industrial test case that is simulation of the
Pellet-Cladding mechanical Interaction (PCI) in Pressurized Water
Reactors (PWR) \cite{Michel13}. Precise simulations of this problem
require cells of 1 $\mu$m for a structure of 1 cm. The LDC multigrid
method could be efficient to overcome this issue. We present results
obtained on the 2D linear elasticity test cases with discontinuous
boundary condition. Whatever the thresholds set by user, the proposed
coupled method gives really satisfactory results. Memory space and CPU
time saving will be also pointed out.

\begin{thebibliography}{1}

\bibitem{Barbie13CAS}
L.~Barbi{\'e}, I.~Rami{\`e}re, and F.~Lebon.
\newblock Strategies around the local defect correction multi-level refinement
method for three-dimensional linear elastic problems.
\newblock {\em Computers and Structures}, 130:73--90, 2014.

\bibitem{Barbie11}
L.~Barbi{\'e}, I.~Rami{\`e}re, F.~Lebon, and J.~Sercombe.
\newblock {AMR} methods in solids mechanics for pellet-cladding interaction
modelling.
\newblock In D.~Aubry, P.~D\'iez, B.~Tie, and N.~Par{\'e}s, editors, {\em V
International Conference on Adaptive Modeling and Simulation}, pages 70--81,
2011.
\newblock 6-8 June, Paris, France.

\bibitem{Hackbusch84}
W.~Hackbusch.
\newblock {Local Defect Correction Method and Domain Decomposition Techniques}.
\newblock {\em Computing Suppl. Springer-Verlag}, 5:89--113, 1984.

\bibitem{Michel13}
B.~Michel, C.~Nonon, J.~Sercombe, F.~Michel, and V.~Marelle.
\newblock Simulation of pellet-cladding interaction with the pleiades fuel
performance software environment.
\newblock {\em Nuclear Technology}, 182, 2013.

\bibitem{Zienkiewicz87}
O.C. Zienkiewicz and J.Z. Zhu.
\newblock A simple error estimator and adaptive procedure for practical
engineering analysis.
\newblock {\em International Journal for Numerical Methods in Engineering},
24:337--357, 1987.

\end{thebibliography}


\end{document}
